\scanpage{114}
% archive foliation: 113

La fonction $\wp$, considérée comme morphisme
\[
(52)\qquad\wp:E\longrightarrow\mathbb P^1_k
\]
rigidise l'iso
\[
(53)\qquad\unsure{\dot\wp}:E/\{\pm\mathrm{id}_E\}\simeq
\text{droite projective type }\mathbb P^1_k
\]
par la condition
\[
(54)\qquad\wp(0_E)=\infty
\]
ce qui détermine $\wp$ modulo transformation affine
\[
(55)\qquad\wp\longmapsto\wp_1=a\wp+b\qquad a\in k^*,\ b\in k.
\]
Les points d'ordre 2
\[
(56)\qquad(\unsure{\varepsilon_i|e_i})_{i\in J},\qquad\unsure{\varepsilon_i|e_i}\in{}_2E(k)^*
\]
sont transformés par $\wp$ en des points
\[
(57)\qquad\unsure{\xi|\zeta}_i=\wp(\unsure{\varepsilon_i|e_i})\in\mathbb P^1_k(k)\setminus\{\infty\}=k
\]
% margin position: left beside (57)
\margin{\uline{NB} On obtient \unsure{ainsi des} $E,\wp$
\unsure{pr} n'importe quel triplet de points distincts de $k$.}
Ce sont, bien sûr, avec le point $\infty$, les points de ramification de $\wp$.
La normalisation de Weierstrass, possible en car. $\ne3$, \deleted{qui} est
\[
(58)\qquad\underbrace{\sum\unsure{\xi|\zeta}_i}_{-a_1}=0
\]
elle détermine $\wp$ modulo transformation
\[
\wp\longmapsto\wp_1=a\wp\qquad a\in k^*.
\]
En général, on pose
\[
(59)\quad\left\{\begin{aligned}
a_1&=-\sum\unsure{\xi|\zeta}_i\\
a_2&=\sum\unsure{\xi|\zeta}_i\unsure{\xi|\zeta}_j\\
a_3&=\deleted{\unsure{{}+\prod\unsure{\xi|\zeta}_i}}-\prod\unsure{\xi|\zeta}_i
\end{aligned}\right.
\]
